\subsection{预解拓扑与函数演算的连续性}

在本节中，E 表示一个复希尔伯特空间。回顾我们用 \(\mathcal{A}(\mathrm{E})\) 表示 E 上自伴部分算子所成的集。我们将把 IV, p. 194 第 8 目的连续性性质推广到 \(\mathcal{A}(\mathrm{E})\) 上。

定义 6. --- 设 T 是 \(\mathcal{L}(\mathrm{E})\) 上的一个局部凸拓扑。\(\mathcal{A}(\mathrm{E})\) 上的 T-预解拓扑是使映射 \(u \mapsto \mathrm{R}(u, \lambda)\) （从 \(\mathcal{A}(\mathrm{E})\) 到赋以拓扑 T 的 \(\mathcal{L}(\mathrm{E})\) 中）对每个 \(\lambda \in \mathbf{C} - \mathbf{R}\) 都连续的最粗拓扑。

命题 11. --- 设 T 是 \(\mathcal{L}(\mathrm{E})\) 上的局部凸拓扑，它比 \(\mathcal{L}(\mathrm{E})\) 的巴拿赫空间拓扑粗。

\begin{enumerate}
\def\labelenumi{\alph{enumi})}
\item
  设 \(f \in \mathcal{C}_0(\mathbf{R})\) 。对每个序列 \((u_n)_{n \in \mathbf{N}}\) （属于 \(\mathcal{A}(\mathrm{E})\) ，按 \(\mathcal{T}\)-预解拓扑收敛于 u），序列 \((f(u_n))_{n \in \mathbf{N}}\) 收敛于 \(f(u)\) ，在空间 \(\mathcal{L}(\mathrm{E})\) 中，该空间赋以拓扑 \(\mathcal{T}\) ；
\item
  假设每个 \(u \in \mathcal{A}(\mathrm{E})\) 对 \(\mathcal{T}\)-预解拓扑都具有可数的基本邻域系。从 \(\mathcal{C}_0(\mathbf{R}) \times \mathcal{A}(\mathrm{E})\) 到空间 \(\mathcal{L}(\mathrm{E})\) （赋以拓扑 \(\mathcal{T}\) ）中的、由 \((f, u) \mapsto f(u)\) 所定义的映射是连续的。
\end{enumerate}

我们来证明 \(a\) )。拓扑 \(\mathcal{T}\) 在 \(\mathcal{L}(\mathrm{E})\) 上是由对拓扑 \(\mathcal{T}\) 连续的半范数所定义的诸拓扑的上确界（EVT, II, p. 26, 推论及其后的注）。因此，只须证明：对每个半范数 \(p\) ，只要它在 \(\mathcal{L}(\mathrm{E})\) 上对拓扑 \(\mathcal{T}\) 连续，序列 \(p(f(u_n) - f(u))\) 就收敛于 0（TG, I, p. 51, 命题 10）。

设 p 是这样一个半范数。用 \(\mathcal{L}(\mathrm{E})_{p}\) 表示赋以半范数 p 的空间 \(\mathcal{L}(\mathrm{E})\) 的分离巴拿赫空间完备化，并用 \(\varpi\) 表示从 \(\mathcal{L}(\mathrm{E})\) 到 \(\mathcal{L}(\mathrm{E})_{p}\) 中的典范映射；它是连续的，且有 \(p(u) = \|\varpi(u)\|_{p}\) ，对所有 \(u \in \mathcal{L}(\mathrm{E})\) ，其中范数是空间 \(\mathcal{L}(\mathrm{E})_{p}\) 的范数。

由于 T 比 \(\mathcal{L}(\mathrm{E})\) 的巴拿赫空间拓扑粗，存在 \(c \geqslant 0\) 使得 \(p(u) \leqslant c \|u\|\) 对所有 \(u \in \mathcal{L}(\mathrm{E})\) 成立。

先假设 \(f \in \mathcal{D}(\mathbf{R})\) 。设 \(\tilde{f}\) 是 \(f\) 的一个几乎解析延拓。用 \(\mu\) 表示 \(\mathbf{C}\) 上的勒贝格测度。由埃尔弗--舍斯特兰德公式（IV, p. 284 的定理 2），有

\[
\varpi (f (u _ {n})) = - \frac {1}{\pi} \int_ {\mathbf {C}} \frac {\partial \widetilde {f}}{\partial \bar {z}} (\lambda) \varpi (\mathrm{R} (u _ {n}, \lambda)) d \mu (\lambda)\tag{14}
\]

对所有 \(n \in \mathbf{N}\) 成立。

设 \(\lambda \in \mathbf{C} - \mathbf{R}\) 。依据假设，预解式序列 \(\mathrm{R}(u_n,\lambda)\) 收敛于 \(\mathrm{R}(u,\lambda)\) ，在空间 \(\mathcal{L}(\mathrm{E})\) 中，该空间赋以拓扑 \(\mathcal{T}\) ，因此序列 \((\varpi (\mathrm{R}(u_n,\lambda)))_n\) 收敛于 \(\varpi (\mathrm{R}(u,\lambda))\) ，在巴拿赫空间 \(\mathcal{L}(\mathrm{E})_p\) 中。

对每个 \(n \in \mathbf{N}\) ，我们有

\[
\left\| \frac {\partial \widetilde {f}}{\partial \bar {z}} (\lambda) \mathrm{R} (u _ {n}, \lambda) \right\| \leqslant \left| \frac {1}{\mathcal{I} (\lambda)} \frac {\partial \widetilde {f}}{\partial \bar {z}} (\lambda) \right|
\]

（IV, p. 248 的命题 17），因此

\[
\left\| \frac {\partial \widetilde {f}}{\partial \overline {{z}}} (\lambda) \varpi (\mathrm{R} (u _ {n}, \lambda)) \right\| _ {p} \leqslant c \left| \frac {1}{\mathcal {I} (\lambda)} \frac {\partial \widetilde {f}}{\partial \overline {{z}}} (\lambda) \right|.
\]

IV, p. 281 的估计式 (13) 表明，这个不等式的右端对 \(\lambda \in \mathbf{C} - \mathbf{R}\) 是有界函数；它在 \(\mathbf{C}\) 上可积，因为 \(\widetilde{f}\) 具有紧支撑。由勒贝格定理（INT, IV, p. 137, § 3, 第 7 目, 定理 6）以及应用于 \(f\) 的埃尔弗--舍斯特兰德公式，我们推出 \(\varpi(f(u_n))\) 收敛于 \(\varpi(f(u))\) ，从而 \(p(f(u_n) - f(u))\) 趋于 0。

现在假设 \(f \in \mathcal{C}_0(\mathbf{R})\) 。设 \(\varepsilon > 0\) 。存在函数 \(f_{\varepsilon}\) ，它属于 \(\mathcal{D}(\mathbf{R})\) ，使得 \(\| f_{\varepsilon} - f \|_{\infty} \leqslant \varepsilon\) （参见 IV, p. 202 的命题 4, a) 及 INT, III, p. 45, § 1, 第 2 目, 命题 3）。于是，依据 IV, p. 275 的命题 5, c)，有 \(\| f(u) - f_{\varepsilon}(u) \| \leqslant \varepsilon\) 及 \(\| f(u_n) - f_{\varepsilon}(u_n) \| \leqslant \varepsilon\) ，对所有 \(n \in \mathbf{N}\) 。因此，可得

\[
p (f (u _ {n}) - f (u)) \leqslant 2 c \varepsilon + p (f _ {\varepsilon} (u _ {n}) - f _ {\varepsilon} (u))
\]

对所有 \(n \in \mathbf{N}\) 成立。由于 \(f_{\varepsilon} \in \mathcal{D}(\mathbf{R})\) ，由前一种情形，序列 \((p(f_{\varepsilon}(u_n) - f_{\varepsilon}(u)))_{n \in \mathbf{N}}\) 收敛于 0，因此 \(p(f(u_n) - f(u))\) 收敛于 0。这就完成了 a) 的证明。

我们来证明断言 \(b\) )。在关于 \(\mathcal{T}\) 的假设下，由 TG, IX, p. 17 的命题 10 及其后的注，以及断言 \(a\) )，可得映射 \(u \mapsto f(u)\) 从 \(\mathcal{A}(\mathrm{E})\) 到 \(\mathcal{L}(\mathrm{E})\) 中是连续的，当 \(f \in \mathcal{C}_0(\mathbf{R})\) 时。由于映射 \(f \mapsto f(u)\) 从 \(\mathcal{C}_0(\mathbf{R})\) 到 \(\mathcal{L}(\mathrm{E})\) 中连续且范数 \(\leqslant 1\) ，对所有 \(u \in \mathcal{A}(\mathrm{E})\) （IV, p. 275 的命题 5, \(c\) )），断言 \(b\) ) 由 TG, X, p. 13, 推论 3 得出。

例. --- 1) 设 \(\mathscr{S}\) 是 E 的有界子集所成的一个集。\(\mathscr{S}\)-拓扑在 \(\mathcal{L}(\mathrm{E})\) 上（EVT, III, p. 13）是比 \(\mathcal{L}(\mathrm{E})\) 的巴拿赫空间拓扑更粗的局部凸拓扑。

\begin{enumerate}
\def\labelenumi{\arabic{enumi})}
\setcounter{enumi}{1}
\tightlist
\item
  设 \(T_{b}\) 是 \(\mathcal{L}(\mathrm{E})\) 的巴拿赫空间拓扑。对 \(T_{b}\)-预解拓扑，每个 \(u \in \mathcal{A}(\mathrm{E})\) 都具有可数的基本邻域系。
\end{enumerate}

事实上，只须证明：对每个 \(\lambda \in \mathbf{C} - \mathbf{R}\) 及每个 \(\varepsilon > 0\) ，存在 \(\lambda' \in \mathbf{Q} + i\mathbf{Q}^*\) 及整数 \(n \geqslant 1\) ，使得每个部分算子 \(v \in \mathcal{A}(\mathrm{E})\) 若满足 \(\| \mathrm{R}(v, \lambda') - \mathrm{R}(u, \lambda')\| < 1/n\) ，则也满足条件 \(\|R(v,\lambda)-R(u,\lambda)\|<\varepsilon\) ；写出

\[
\begin{array}{c} \| \mathrm{R} (v, \lambda) - \mathrm{R} (u, \lambda) \| \leqslant \| \mathrm{R} (v, \lambda) - \mathrm{R} (v, \lambda^ {\prime}) \| \\ + \| \mathrm{R} (v, \lambda^ {\prime}) - \mathrm{R} (u, \lambda^ {\prime}) \| + \| \mathrm{R} (u, \lambda^ {\prime}) - \mathrm{R} (u, \lambda) \|, \end{array}
\]

这个性质由 IV, p. 245 的公式 (8)、IV, p. 248 的命题 17 的诸估计式，以及 \(\mathbf{Q} + i\mathbf{Q}^*\) 在 \(\mathbf{C} - \mathbf{R}\) 中处处稠密这一事实得出。

若把 \(\mathcal{C}_{0}(\mathbf{R})\) 换成 \(\mathcal{C}_{b}(\mathbf{R})\) ，命题的结论一般不成立（IV, p. 366 习题 29, e)）。尽管如此，我们有下列结果。

推论. --- 设 \(\mathcal{T}_{s}\) 是 \(\mathcal{L}(\mathrm{E})\) 上的逐点收敛拓扑。设 \(f \in \mathcal{C}_{b}(\mathbf{R})\) 。对每个序列 \((u_{n})_{n \in \mathbf{N}}\) （属于 \(\mathcal{A}(\mathrm{E})\) ，收敛于 \(u\) ，按 \(\mathcal{T}_{s}\)-预解拓扑），序列 \((f(u_{n}))_{n \in \mathbf{N}}\) 收敛于 \(f(u)\) ，在 \(\mathcal{L}(\mathrm{E})\) 中，该空间赋以拓扑 \(\mathcal{T}_{s}\) 。

设 \((u_{n})\) 是 \(\mathcal{A}(\mathrm{E})\) 中按 \(T_{s}\)-预解拓扑收敛于 u 的序列。设 \(x \in E\) ，并设 \(\varepsilon > 0\) 。

对每个整数 \(N \in \mathbf{N}\) ，设 \(\varphi_{N} \in \mathcal{K}(\mathbf{R})\) 是支撑包含在 \([-(N+1), N+1]\) 中的函数，使得 \(0 \leqslant \varphi_{N} \leqslant 1\) ，且 \(\varphi_{N}(t) = 1\) 对所有 \(t \in [-N, N]\) 成立。当 N 趋于无穷时，函数 \(\varphi_{N}\) 逐点收敛于 1，且满足 \(|\varphi_{N}| \leqslant 1\) ；因此 IV, p. 276 的命题 6 蕴含存在 \(N \in \mathbf{N}\) 使得 \(\|\varphi_{N}(u)x - x\| \leqslant \varepsilon\) 。置 \(f_{N} = f\varphi_{N}\) 。

对每个 \(n\in \mathbf{N}\) ，我们有

\[
\begin{array}{r l} \| f (u _ {n}) x - f (u) x \| \leqslant & \| f (u _ {n}) x - f _ {\mathrm{N}} (u _ {n}) x \| + \\ & \| f _ {\mathrm{N}} (u _ {n}) x - f _ {\mathrm{N}} (u) x \| + \| f _ {\mathrm{N}} (u) x - f (u) x \|. \end{array} \tag {15}
\]

对每个 \(v\in \mathcal{A}(\mathrm{E})\) ，我们有

\[
\begin{array}{r l} & {\| f (v) x - f _ {\mathrm{N}} (v) x \| = \| (f (1 - \varphi_ {\mathrm{N}})) (v) x \|} \\ & {\qquad \leqslant \| f (v) \|   \| x - \varphi_ {\mathrm{N}} (v) x \| \leqslant \| f \| _ {\infty}   \| x - \varphi_ {\mathrm{N}} (v) x \|} \end{array}
\]

（IV, p. 275 的命题 5, c)）。由于 \(\varphi_{\mathrm{N}} \in \mathcal{C}_0(\mathbf{R})\) ，命题 11 的断言 \(a\) ) 应用于 \(\mathcal{T}_s\)-预解拓扑，蕴含 \(\varphi_{\mathrm{N}}(u_n)x\) 收敛于 \(\varphi_{\mathrm{N}}(u)x\) ，当 \(n \to +\infty\) 时。因此，对所有充分大的 \(n\) ，有

\[
\left\| x - \varphi_ {\mathrm{N}} (u _ {n}) x \right\| \leqslant \left\| x - \varphi_ {\mathrm{N}} (u) x \right\| + \varepsilon \| x \| \leqslant (1 + \| x \|) \varepsilon .
\]

于是不等式 (15) 变成

\[
\| f (u _ {n}) x - f (u) x \| \leqslant \| f \| _ {\infty} (2 + \| x \|) \varepsilon + \| f _ {\mathrm{N}} (u _ {n}) x - f _ {\mathrm{N}} (u) x \|
\]

对所有充分大的 n 成立。对函数 \(f_{\mathrm{N}} \in \mathcal{C}_{0}(\mathbf{R})\) 及 \(T_{s}\)-预解拓扑应用 loc. cit. 的断言 a)，即完成证明。

